\documentclass{article}
\usepackage{amsthm}
\newtheorem{theorem}{Theorem}
\newtheorem{lemma}[theorem]{Lemma}
\newtheorem{corollary}[theorem]{Corollary}
\begin{document}
\title{A paper with an ordinary shape}
\maketitle

\begin{abstract}
We prove a bound; see Theorem~\ref{thm:main} for the statement.
\end{abstract}

\input{intro}

\section{Preliminaries}

\begin{lemma}\label{lem:key}
Every widget is a gadget.
\end{lemma}

\begin{theorem}[Main Theorem]\label{thm:main}
Our main result is that every gadget of order $n$ admits a decomposition into
at most $n/2$ widgets, and moreover the decomposition can be computed in
polynomial time. This statement is deliberately long so that it reads as a
substantial statement rather than a one-line remark, which is one of the
signals the advisory ranking uses when it orders the candidate list.
\end{theorem}

\begin{corollary}
Every gadget of order two is a widget.
\end{corollary}

\end{document}
